\section{Conclusiones}
\label{sec:conclusiones}

\begin{itemize}
    \item El requerimiento de telemetría (16~kb/s útiles con margen 2:1) se cubrió con BPSK a $R_s=32$~ksym/s, SPS~$=8$ y conformación RRC con $\beta=0.35$. El ancho de banda ocupado bajó de unos 64~kHz con pulso rectangular a 43.2~kHz, sin interferencia intersimbólica en los instantes de decisión.
    \item El ruido impulsivo periódico se identifica en el dominio del tiempo y en la constelación, pero casi no en la PSD. Su potencia media es pequeña ($A_{imp}^2/256$) y solo se manifiesta como un aumento del piso y un peine de líneas a la frecuencia de repetición de 1~kHz. Para este tipo de perturbación, la PSD no es un buen indicador del riesgo de error.
    \item El filtro acoplado fue la principal defensa contra los impulsos: en los instantes de decisión reduce un impulso de una muestra a $0.100\,A_{imp}$. Con AWGN de amplitud 0.10, la BER fue cero en todas las pruebas ($A_{imp}$ hasta 5), aunque el margen de decisión bajó de 0.88 a 0.40. En la fase de canal, en cambio, con pulso rectangular y sin filtro acoplado, $A_{imp}=1$ bastaba para llevar símbolos al umbral.
    \item El instante de muestreo resultó crítico: un error de medio símbolo eleva la BER a 0.25 aunque el canal no cambie. Calcular el retardo total del enlace (88 muestras) y, en hardware, recuperar el reloj y calibrar el retardo de bits son condiciones necesarias para que la medición de BER tenga sentido.
    \item Sin impulsos, el enlace alcanza BER $=10^{-3}$ con $E_b/N_0=7.3$~dB, de acuerdo con la teoría de BPSK con codificación diferencial; esta cuesta unos 0.5~dB frente a BPSK coherente, a cambio de eliminar la ambigüedad de fase de 180° del receptor de hardware. Con $A_{imp}=1$ no hay penalización medible. Con $A_{imp}=5$ la penalización es de 1.8~dB y la curva decae lentamente, porque uno de cada 32 símbolos pierde la mitad de su distancia al umbral. Si la amplitud de los impulsos de la planta fuera de ese orden, convendría agregar codificación de canal con entrelazado o limitar la amplitud de las muestras antes del filtro acoplado.
    \item La migración a hardware se hizo sin cambiar ningún parámetro de diseño: modulación, $R_s$, SPS, $\beta$, span, mapeo y decisión se conservaron, y solo se agregaron la interfaz del radio (frecuencia central, tasa del radio, remuestreo, ganancias y desplazamiento del oscilador) y la sincronización de frecuencia, reloj y fase que exigen dos osciladores independientes. La codificación diferencial, cuyo costo medido fue de unos 0.5~dB, es la que permite resolver la ambigüedad de fase del lazo de portadora.
    \item Como trabajo futuro queda medir la curva BER del enlace BPSK propio con los radios, conectados por cable con atenuación, y compararla con la referencia simulada usando el criterio de penalización máxima de 2~dB en BER $=10^{-3}$. También queda evaluar codificación de canal con entrelazado y limitación de amplitud antes del filtro acoplado frente a impulsos de mayor amplitud.
\end{itemize}
